%%%%%%%%%%%%%%%%%%%%%%%%%%%%%%%%%%%%%%%%%%%%%%%%%
\section{重正化}
\label{sec:renormalization}

本节针对上一节辨认出的、可能给夸克准 PDF 带来紫外发散贡献的三类拓扑图，给出其重正化的一般论证。

首先考虑图~\ref{fig:multiloopdiv}(a) 一般图的幂次紫外发散；它可以表示为由 1PI 图构成的诸图之和，如图~\ref{fig:sumSE} 所示。由于紫外发散在坐标空间局域，我们给图~\ref{fig:sumSE} 中第``i''个圆团指定一个特定点 $r_i$。把图~\ref{fig:sumSE} 中``第 1 个''圆团的线性紫外幂次发散表示为 $c \int_0^{r} d r_1$，其中微扰系数 $c=c^{(1)}+c^{(2)}+\cdots$，而 $c^{(1)}=-\frac{\alpha_s C_F}{\pi}\frac{1}{a} $ 是按 $\alpha_s$ 展开的首阶项，如式~(\ref{eq:1a}) 所导出。

借助图~\ref{fig:sumSE} 所示 1PI 图的几何级数求和，把位于 0 与 $r\,(r>0)$ 之间的各 1PI 图的贡献对所有阶求和，可以导出图~\ref{fig:multiloopdiv}(a) 一般图的全部线性紫外幂次发散贡献：
\begin{eqnarray}
&\ & 1+ c \int_0^{r} d r_1 + c^2 \int_0^{r} d r_1 \int_{r_1}^{r} d r_2 + \cdots
\nonumber \\
&\ & \hskip 0.3in
= {\cal P} e^{\, c \int_0^r dr'} = e^{\, c\, r}\, ,
\end{eqnarray}
其中 ${\cal P}$ 表示路径排序。
因此，可以引入一个整体因子 $e^{-c\, |\xi_z|}$ 来消除夸克准 PDF 的所有线性幂次紫外发散。这个整体因子可以理解为沿规范链接运动的试验粒子的质量重正化。以这种方式重正化幂次发散发散最早由文献 \cite{Polyakov:1980ca} 提出。

除幂次紫外发散外，图~\ref{fig:multiloopdiv}(a) 还有对数紫外发散。众所周知 \cite{Dotsenko:1979wb}，这些发散可以通过试验粒子的``波函数''重正化 $Z_{wq}^{-1}$ 消除。

图~\ref{fig:multiloopdiv}(b) 的一般图只有对数紫外发散，实质上来自胶子与规范链接耦合的圈修正。文献 \cite{Dotsenko:1979wb} 已证明这些对数紫外发散可以被 QCD 的耦合常数重正化吸收。因此，若使用重正化的 QCD 拉氏量，则无需为此担心。

最后考察图~\ref{fig:multiloopdiv}(c) 一般图的紫外发散。与图~\ref{fig:multiloopdiv}(a)、(b) 的一般图不同，图~\ref{fig:multiloopdiv}(c) 的圈动量穿过活跃夸克（对胶子准 PDF 的情形则为活跃胶子）。由于图~\ref{fig:multiloopdiv}(c) 各图的紫外发散实质上来自夸克--规范链接顶点的高阶圈修正，而这并非 QCD 拉氏量中的基本耦合，因此使用重正化的 QCD 拉氏量进行计算无助于消除这类紫外发散。也就是说，如果想要可重正的准 PDF，就必须对定义准 PDF 的算符进行重正化以消除这类紫外发散。

那么关键问题就是：定义准 PDF 的算符在重正化下会不会与其他算符混合？如果会，就存在这样的危险：定义准 PDF 的算符在重正化下可能不构成封闭集合。

最低阶的夸克--规范链接顶点，对夸克准 PDF 而言是一个简单的伽马矩阵 $\gamma_z$（对 $A_z=0$ 情形亦相同）。正如上一节所证明的，紫外发散只来自所有圈动量都非常大的区域，因此如果只关心领头的紫外发散，可以令 $p=0$；而如第 \ref{sec:powercounting} 节所示，该领头发散是对数型的。此外，对数紫外发散（$a\to 0$ 时的 $\ln(a)$）在坐标空间局域，不直接依赖 $\xi_z$。这就是说，我们发现图~\ref{fig:multiloopdiv}(c) 的紫外发散项只依赖矢量 $n_z^\mu$，因而以一个常数系数正比于 $\gamma_z$，亦即正比于最低阶的夸克--规范链接顶点。因此，一个常数抵消项就足以消除这类紫外发散。运用簿记式的森林减除方法，可以直接消去高阶发散并证明其净效果是为夸克--规范链接顶点引入一个常数乘法重正化因子 $Z_{vq}^{-1}$。

总之，在费曼规范下使用重正化的 QCD 拉氏量，我们发现夸克准 PDF 余下的全部微扰紫外发散都可以通过引入乘法重正化因子来消除，该乘法重正化因子可在 QCD 微扰论中逐阶计算。我们预期类似的论证也应适用于胶子准 PDF。于是，可以把重正化的坐标空间准 PDF 定义为
\begin{align}\label{eq:renor}
\begin{split}
&\tilde{F}^{R}_{i/p}(\xi_z,\tilde{\mu}^2,p_z)=e^{-C_i |\xi_z|}Z_{wi}^{-1}Z_{vi}^{-1}\tilde{F}^{b}_{i/p}(\xi_z,\tilde{\mu}^2,p_z),
\end{split}
\end{align}
其中 $C_i$、$Z_{wi}$ 与 $Z_{vi}$ 是依赖部分子味``$i$''但不依赖 $\xi_z$ 的重正化常数，并可按 $\alpha_s$ 幂次在 QCD 微扰论中逐阶微扰计算。我们的结论是：坐标空间准 PDF 是可重正的，并且它们在重正化群方程下既不彼此混合，也不与任何其他算符混合。
